\providecommand{\iquest}[1]{¿#1}
\section{Guion de entrevista semiestructurada}
\sloppy

\textbf{C\'odigo del instrumento:} INS-ENT-01. \textbf{Duraci\'on estimada:} 35 a 50 minutos.
\textbf{Dirigido a:} personal de atenci\'on, cocina, caja e inventario de
El Balcón Restaurant. \textbf{Registro:} notas an\'onimas con c\'odigo de rol; sin
grabaci\'on salvo autorizaci\'on expresa.

\textbf{Origen del instrumento.} Este guion reproduce, sin respuestas de ejemplo,
las 45 preguntas de \texttt{guia\_entrevista\_restaurante1\_corregida.pdf},
documento UTEQ empleado en la aplicaci\'on de campo. La versi\'on corregida se
incorpora como instrumento en blanco; las versiones con respuestas redactadas
(`003.pdf`, `004.pdf` y `005.pdf') se usaron solo como apoyo para redactar el
guion y \emph{no} como resultados de la entrevista.

\subsection*{Apertura y consideraciones \'eticas}
Se lee al participante:

\begin{quote}
\small
``Esta entrevista se realiza con el prop\'osito de levantar informaci\'on sobre
la operaci\'on diaria del restaurante: las actividades que realiza el personal,
la forma en que se atienden y registran los pedidos, se controla el inventario,
se gestionan las compras, el personal y la informaci\'on administrativa. La
informaci\'on recolectada permitir\'a describir la situaci\'on actual del
establecimiento y servir\'a \'unicamente como insumo para un trabajo acad\'emico.
La entrevista no eval\'ua el desempe\~no de la persona entrevistada ni busca
proponer soluciones en esta etapa. La participaci\'on es voluntaria; no se
solicitan nombre, c\'edula, tel\'efono ni correo. Puede interrumpir la entrevista
o abstenerse de responder sin justificaci\'on. \iquest{Autoriza el uso an\'onimo
de sus respuestas para este proyecto?}''
\end{quote}

Si se autoriza grabaci\'on, se anota el c\'odigo del participante, la fecha, la
hora de inicio y de finalizaci\'on y la duraci\'on real. El original firmado y la
grabaci\'on se custodian en carpeta privada \texttt{@uteq.edu.ec}.

\subsection*{1. Apertura y perfil del cargo}
\begin{enumerate}[leftmargin=18pt,itemsep=0pt]
    \item \iquest{Cu\'al es su cargo y cu\'anto tiempo lleva en el restaurante?}
    \item \iquest{Cu\'ales son sus responsabilidades en un d\'ia normal y con qu\'e \'areas interact\'ua m\'as?}
    \item \iquest{Qu\'e medio o forma utiliza para registrar actualmente las actividades de su trabajo?}
\end{enumerate}

\subsection*{2. Gesti\'on de pedidos y ventas}
\begin{enumerate}[leftmargin=18pt,itemsep=0pt,resume]
    \item \iquest{C\'omo se atiende un pedido, desde que llega el cliente hasta que paga?}
    \item \iquest{C\'omo se registra cada pedido y cu\'antos se atienden en un d\'ia normal y en uno de mayor movimiento?} D\'ia normal: \rule{1.6cm}{0.4pt}; d\'ia de mayor movimiento: \rule{1.6cm}{0.4pt}.
    \item \iquest{Qu\'e es lo que m\'as se le dificulta cuando hay mucha demanda?}
    \item \iquest{Con qu\'e frecuencia ocurren errores en los pedidos?} a) Nunca; b) Pocas veces al mes; c) Varias veces por semana; d) Todos los d\'ias. \iquest{Por qu\'e cree que ocurren?}
    \item \iquest{C\'omo se realiza el cobro y la facturaci\'on?} \iquest{Qu\'e medios de pago se aceptan y qu\'e inconvenientes ha tenido al cobrar, dividir cuentas o facturar?}
\end{enumerate}

\subsection*{3. Gesti\'on de delivery y entrega a domicilio}
\begin{enumerate}[leftmargin=18pt,itemsep=0pt,resume]
    \item \iquest{El restaurante ofrece servicio de entrega a domicilio?} \iquest{C\'omo se solicita y con qu\'e frecuencia se utiliza?}
    \item \iquest{C\'omo se recibe y registra un pedido para entrega a domicilio?}
    \item \iquest{C\'omo se organiza la preparaci\'on y el despacho de los pedidos a domicilio para que lleguen a tiempo?}
    \item \iquest{Qu\'e medio de transporte o repartidores se utilizan y c\'omo se asigna cada entrega?}
    \item \iquest{Qu\'e problemas ocurren con m\'as frecuencia en las entregas a domicilio (demoras, pedidos incompletos, direcciones, cobros) y c\'omo se resuelven?}
\end{enumerate}

\subsection*{4. Gesti\'on de reservas y atenci\'on al cliente}
\begin{enumerate}[leftmargin=18pt,itemsep=0pt,resume]
    \item \iquest{C\'omo se registra y confirma una reserva?} \iquest{C\'omo se decide qu\'e mesa se asigna a cada cliente o grupo?}
    \item \iquest{Con qu\'e frecuencia hay problemas con mesas o reservas (cruces, esperas, cancelaciones)?} a) Nunca; b) Rara vez; c) Con frecuencia; d) Casi siempre. \iquest{Por qu\'e cree que ocurren?}
    \item \iquest{C\'omo se atienden las solicitudes especiales de los clientes (alergias, celebraciones, preferencias)?}
    \item \iquest{Qu\'e es lo que m\'as reclaman los clientes?} a) Tiempo de espera; b) Errores en el pedido; c) Trato del personal; d) Precios o cobros; e) Otro. \iquest{Por qu\'e cree que ocurre?}
\end{enumerate}

\subsection*{5. Gesti\'on de cocina y preparaci\'on de alimentos}
\begin{enumerate}[leftmargin=18pt,itemsep=0pt,resume]
    \item \iquest{C\'omo llega la comanda a la cocina y c\'omo se decide el orden en que se preparan los pedidos?}
    \item \iquest{C\'omo se comunican cocina y meseros cuando un pedido est\'a listo o cambia?}
    \item \iquest{Con qu\'e frecuencia se retrasan los platos?} a) Nunca; b) Algunas veces; c) Frecuentemente; d) Casi siempre. \iquest{Cu\'al es la causa principal?}
    \item \iquest{C\'omo se controla que los platos se preparen con la misma receta y porci\'on?} \iquest{Qu\'e hacen cuando falta un ingrediente en pleno servicio?}
\end{enumerate}

\subsection*{6. Gesti\'on de inventario}
\begin{enumerate}[leftmargin=18pt,itemsep=0pt,resume]
    \item \iquest{C\'omo se controlan los productos e insumos y c\'omo se registran las entradas y salidas?}
    \item \iquest{Cada cu\'anto se revisa el inventario?} a) Diariamente; b) Semanalmente; c) Mensualmente; d) Solo cuando falta algo. \iquest{Por qu\'e con esa frecuencia?}
    \item \iquest{C\'omo y d\'onde se almacenan los insumos, sobre todo los perecederos?}
    \item \iquest{Con qu\'e frecuencia hay productos vencidos, da\~nados o desperdiciados?} \iquest{C\'omo se decide cu\'ando y cu\'anto reponer de cada producto?}
\end{enumerate}

\subsection*{7. Gesti\'on de compras y proveedores}
\begin{enumerate}[leftmargin=18pt,itemsep=0pt,resume]
    \item \iquest{C\'omo se solicitan los productos a los proveedores y qui\'en decide qu\'e se compra y en qu\'e cantidades?}
    \item \iquest{Con cu\'antos proveedores trabajan y c\'omo se verifica que lo recibido coincida con lo pedido y lo facturado?}
    \item \iquest{Con qu\'e frecuencia hay problemas con proveedores (retrasos, faltantes, calidad)?} a) Nunca; b) Rara vez; c) Con frecuencia; d) Casi siempre. \iquest{Cu\'ales son los problemas m\'as comunes y por qu\'e cree que ocurren?}
    \item \iquest{C\'omo se registran las compras y los pagos a proveedores?} \iquest{C\'omo se elige entre proveedores o se comparan precios?}
\end{enumerate}

\subsection*{8. Gesti\'on del personal}
\begin{enumerate}[leftmargin=18pt,itemsep=0pt,resume]
    \item \iquest{Cu\'antas personas trabajan en el restaurante y en qu\'e \'areas?} Total: \rule{1.5cm}{0.4pt}; \'areas: \rule{3cm}{0.4pt}.
    \item \iquest{C\'omo se asignan las tareas a cada empleado y c\'omo se organizan los turnos?}
    \item \iquest{Con qu\'e frecuencia hay ausencias o cambios de \'ultimo minuto?} a) Nunca; b) Rara vez; c) Con frecuencia; d) Casi siempre. \iquest{Por qu\'e cree que ocurre y qu\'e se hace cuando pasa?}
    \item \iquest{C\'omo se coordinan las actividades entre cocina, meseros y caja?} \iquest{C\'omo se supervisa o eval\'ua el desempe\~no del personal?}
\end{enumerate}

\subsection*{9. Gesti\'on administrativa y reportes}
\begin{enumerate}[leftmargin=18pt,itemsep=0pt,resume]
    \item \iquest{C\'omo se registran los ingresos y gastos?}
    \item \iquest{Qu\'e informes elabora o recibe actualmente?} \iquest{Cada cu\'anto y qui\'en los usa?}
    \item \iquest{C\'omo analiza qu\'e productos se venden m\'as o menos?} \iquest{Qu\'e tan f\'acil es obtener la informaci\'on que necesita para decidir?}
\end{enumerate}

\subsection*{10. Dificultades del trabajo diario}
\begin{enumerate}[leftmargin=18pt,itemsep=0pt,resume]
    \item \iquest{Qu\'e tareas de su trabajo le consumen m\'as tiempo o esfuerzo?}
    \item \iquest{Qu\'e informaci\'on del restaurante es m\'as dif\'icil de obtener o mantener actualizada?}
    \item \iquest{Qu\'e ocurre cuando la persona que conoce un proceso no est\'a presente?}
\end{enumerate}

\subsection*{11. Conocimiento y expectativas sobre inteligencia artificial}
\begin{enumerate}[leftmargin=18pt,itemsep=0pt,resume]
    \item \iquest{Qu\'e entiende usted por inteligencia artificial y qu\'e tanto conoce sobre ella?} a) Nada; b) He escuchado el t\'ermino; c) Entiendo lo b\'asico; d) La uso con frecuencia.
    \item \iquest{Ha usado alguna herramienta con inteligencia artificial, en el trabajo o en su vida personal?} \iquest{Cu\'al y c\'omo fue su experiencia?}
    \item \iquest{Qu\'e opina de que un restaurante use inteligencia artificial en su operaci\'on y qu\'e tanto confiar\'ia en su uso?} \iquest{Por qu\'e?}
    \item \iquest{Qu\'e es lo que m\'as le preocupar\'ia sobre el uso de la inteligencia artificial en el trabajo?} a) Que cometa errores; b) Perder el control de las decisiones; c) Que sea dif\'icil de aprender; d) La seguridad de los datos; e) Que reemplace funciones del personal; f) Nada me preocupa; g) Otro. \iquest{Por qu\'e eligi\'o esa opci\'on?}
    \item \iquest{Qu\'e informaci\'on del restaurante considera que deber\'ia mantenerse bajo control humano y no ser manejada por inteligencia artificial?} \iquest{Por qu\'e?}
\end{enumerate}

\subsection*{12. Cierre}
\begin{enumerate}[leftmargin=18pt,itemsep=0pt,resume]
    \item \iquest{Hay alg\'un proceso o problema importante del que no hayamos hablado o algo que desee agregar?}
\end{enumerate}

\noindent\textbf{Observaciones del entrevistador:}\\[0.2cm]
\rule{\linewidth}{0.4pt}\\[0.6cm]
\rule{\linewidth}{0.4pt}
